% Vietnamese translation for OI Public Library
% Source-File: IOI中国国家候选队论文/2007/day1/4.余江伟《如何解决好动态统计问题》.ppt
\documentclass[11pt]{article}
\usepackage{oipl-vn}

\newcommand{\Oh}{\mathcal{O}}

\title{Bản trình chiếu: Cách giải tốt các bài toán thống kê động}
\author{Yu Jiangwei}
\date{2007}

\begin{document}
\maketitle

\origtitle{Slide companion for dynamic statistics problems}
\sourcepath{IOI China candidate team papers / 2007 / day 1 / Yu Jiangwei.ppt}

\section{Phạm vi bản dịch}

Tệp gốc chỉ có PowerPoint. Phần chữ đọc được cho thấy bài trình bày xoay quanh
các bài thống kê động trên dãy và cây tìm kiếm cân bằng, đặc biệt là splay.

\section{Bối cảnh bài toán}

Các bài thống kê động thường yêu cầu vừa cập nhật dữ liệu vừa trả lời truy vấn
thứ hạng, tổng, cực trị hoặc đoạn tối ưu. Nếu làm lại từ đầu sau mỗi thao tác,
chi phí dễ tăng tới tuyến tính hoặc cao hơn cho mỗi truy vấn. Vì vậy cần một
cấu trúc dữ liệu duy trì đủ thông tin phụ trong khi vẫn hỗ trợ thay đổi cục
bộ.

Một ví dụ được slide nhắc tới là dạng bài quản lý dãy với số phần tử có thể
lên tới \(10^6\), cùng các lệnh như \texttt{INSERT}, \texttt{DELETE},
\texttt{MAKE-SAME}, \texttt{REVERSE}, \texttt{GET-SUM} và
\texttt{MAX-SUM}. Đây là kiểu yêu cầu cần cây biểu diễn dãy, không chỉ cây
tìm kiếm theo khóa.

\section{Vai trò của splay}

Splay tree đưa nút vừa truy cập lên gốc bằng các bước Zig, Zag và các tổ hợp
hai bước. Tính chất quan trọng là chi phí khấu hao \(\Oh(\log n)\), trong khi
cài đặt có thể xử lý tách, ghép và truy cập phần tử thứ \(k\) khá trực tiếp.

Khi dùng splay để biểu diễn dãy, khóa ngầm là vị trí trong dãy. Mỗi nút lưu
kích thước cây con, tổng đoạn, tổng tiền tố tốt nhất, tổng hậu tố tốt nhất và
tổng đoạn con tốt nhất. Cờ đảo đoạn hoặc gán cùng giá trị được đẩy xuống khi
cần, nhờ đó các thao tác trên đoạn liên tục trở thành tách đoạn ra, chỉnh cờ
hoặc thông tin, rồi ghép lại.

\section{Ghi chú triển khai}

Điểm dễ sai nằm ở cập nhật thông tin sau phép quay và khi đẩy cờ lười. Hàm
\textsc{SplayKth} phải tìm đúng phần tử thứ \(k\) theo kích thước cây con;
sau đó có thể đưa hai biên của đoạn cần xử lý vào hai vị trí cố định để cô
lập đoạn ở một cây con. Cách tổ chức này là khuôn mẫu cho nhiều bài thống kê
động có thao tác trên đoạn.

\end{document}
